\begin{abstract}
With known gain, an ordered orthogonal scan minimizes conditioning-driven amplification of additive noise; with unknown slow multiplicative gain and unconstrained objects, the same deterministic chronology creates an exact gain--object ambiguity. We resolve this paradox with an identifiability theory of blind gain--object separation for single-pixel (ghost) imaging, distinguishing algebraic identifiability, statistical identifiability, and estimation conditioning. Every square invertible design fails algebraically (modulo global scale); tall designs restore uniqueness generically at $N \ge K+p$ frames for a $p$-dimensional drift class. Independently, a stationarity anchor pins the relative gain statistically: once the object carrier is stationary in time, slow modulation of the record can only be gain, recovered at the minimax rate by a windowed log-domain estimator. Acquisition by the randomized Hadamard transform (SRHT, at full sampling) preserves orthogonal conditioning while supplying, under Walsh-flatness or a successful random permutation of a sufficiently spread object, the stationary or exchangeable carrier the estimator needs; an exact finite-noise relative mean-square-error identity then yields a design chart with a computable flip boundary between ordered and randomized acquisition, whose gain-known skeleton is validated on the simulation grid with no free parameters (raw-metric median agreement factor 1.04). In frame-matched simulations, raw ordered Hadamard leaves a non-decaying blind gain error of 0.94--0.99 across all averaging windows, where randomized arms reach 0.007--0.015 (medians with two-way clustered 95\% confidence intervals), and randomized Hadamard acquisition recovers $+5.45$ dB under slow drift. The acquisition rule: randomize the carrier's chronology --- make the object stop pretending to be time.
\end{abstract}

\begin{IEEEkeywords}
ghost imaging; single-pixel imaging; blind calibration; identifiability; bilinear inverse problems
\end{IEEEkeywords}
